\hypertarget{ej10}{}
\noindent \textbf{Ejercicio 10.} Sean $P(x:\mathbb{Z})$ y $Q(x:\mathbb{Z})$ dos predicados cualesquiera que nunca se indefinen. Escribir el predicado asociado a cada uno de los siguientes enunciados:

\begin{enumerate}[label=\alph*)]
    \item “Existe un único número natural menor a 10 que cumple P”
    \item “Existen al menos dos números naturales menores a 10 que cumplen P”
    \item “Existen exactamente dos números naturales menores a 10 que cumplen P”
    \item “Todos los enteros pares que cumplen P, no cumplen Q”
    \item “Si un entero cumple P y es impar, no cumple Q”
    \item “Todos los enteros pares cumplen P, y todos los enteros impares que no cumplen P cumplen Q”
    \item “Si hay un número natural menor a 10 que no cumple P entonces ninguno natural menor a 10 cumple Q; y si todos los naturales menores a 10 cumplen P entonces hay al menos dos naturales menores a 10 que cumplen Q”
\end{enumerate}

\vspace{0.2cm}
{\color{gray}\hrule}
\vspace{0.4cm}

\textbf{Resolución:}

\textit{Nota: Siguiendo la convención de la materia vista en el Ejercicio 9, traducimos "número natural menor a 10" usando el dominio $\mathbb{Z}$ con el rango $0 \leq x < 10$.}

\begin{enumerate}[label=\textbf{\alph*)}]

    \item \textbf{“Existe un único número natural menor a 10 que cumple P“} \\
    Para asegurar que es "único", primero afirmamos que existe al menos uno que cumple. Luego, garantizamos que cualquier otro número que también cumpla, debe ser obligatoriamente el mismo que ya encontramos.
    \[ (\exists x: \mathbb{Z})(0 \leq x < 10 \wedge P(x) \wedge (\forall y: \mathbb{Z})(0 \leq y < 10 \wedge P(y) \rightarrow x = y)) \]

    \vspace{0.4cm}
    \item \textbf{“Existen al menos dos números naturales menores a 10 que cumplen P”} \\
    Pedimos la existencia de dos variables distintas. El "al menos" significa que no nos importa si hay 3, 4 o más; con encontrar dos diferentes es suficiente. La clave es exigir que $x \neq y$.
    \[ (\exists x, y: \mathbb{Z})(0 \leq x < 10 \wedge 0 \leq y < 10 \wedge x \neq y \wedge P(x) \wedge P(y)) \]

    \vspace{0.4cm}
    \item \textbf{“Existen exactamente dos números naturales menores a 10 que cumplen P”} \\
    Es la combinación de las dos lógicas anteriores. Afirmamos que existen dos distintos (al menos dos), y luego usamos el $\forall$ para "cerrar la puerta" y asegurar que cualquier otro elemento $z$ que intente cumplir $P$, tiene que ser idéntico a $x$ o a $y$.
    \begin{multline*}
    (\exists x, y: \mathbb{Z})(0 \leq x < 10 \wedge 0 \leq y < 10 \wedge x \neq y \wedge P(x) \wedge P(y) \wedge \\
    (\forall z: \mathbb{Z})(0 \leq z < 10 \wedge P(z) \rightarrow (z = x \vee z = y)))
    \end{multline*}

    \vspace{0.4cm}
    \item \textbf{“Todos los enteros pares que cumplen P, no cumplen Q”} \\
    La paridad de un entero se expresa analizando su resto en la división por 2 ($x \bmod 2 = 0$). Como habla de "todos", usamos un $\forall$ con una implicación.
    \[ (\forall x: \mathbb{Z})(x \bmod 2 = 0 \wedge P(x) \rightarrow \neg Q(x)) \]

    \vspace{0.4cm}
    \item \textbf{“Si un entero cumple P y es impar, no cumple Q”} \\
    Es un caso análogo al anterior. La frase "si un entero..." es una forma condicional que aplica universalmente. Que sea impar se escribe como $x \bmod 2 \neq 0$.
    \[ (\forall x: \mathbb{Z})(P(x) \wedge x \bmod 2 \neq 0 \rightarrow \neg Q(x)) \]

    \vspace{0.4cm}
    \item \textbf{"Todos los enteros pares cumplen P, y todos los enteros impares que no cumplen P cumplen Q"} \\
    Son dos afirmaciones universales distintas unidas por una gran conjunción. Podemos dividirlo claramente en dos cuantificadores $\forall$ separados.
    \[ (\forall x: \mathbb{Z})(x \bmod 2 = 0 \rightarrow P(x)) \wedge (\forall y: \mathbb{Z})(y \bmod 2 \neq 0 \wedge \neg P(y) \rightarrow Q(y)) \]

    \vspace{0.4cm}
    \item \textbf{“Si hay un número natural menor a 10 que no cumple P entonces ninguno natural menor a 10 cumple Q; y si todos los naturales menores a 10 cumplen P entonces hay al menos dos naturales menores a 10 que cumplen Q”} \\
    Este es un mega-enunciado formado por dos grandes implicaciones unidas por un "y" lógico ($\wedge$). Traducimos cada bloque por separado: \\
    \textit{Bloque 1:} $(\exists x: \mathbb{Z})(0 \leq x < 10 \wedge \neg P(x)) \rightarrow (\forall y: \mathbb{Z})(0 \leq y < 10 \rightarrow \neg Q(y))$ \\
    \textit{Bloque 2:} $(\forall z: \mathbb{Z})(0 \leq z < 10 \rightarrow P(z)) \rightarrow (\exists v, w: \mathbb{Z})(0 \leq v < 10 \wedge 0 \leq w < 10 \wedge v \neq w \wedge Q(v) \wedge Q(w))$ \\
    Juntando todo queda:
    \begin{multline*}
    \Big( (\exists x: \mathbb{Z})(0 \leq x < 10 \wedge \neg P(x)) \rightarrow (\forall y: \mathbb{Z})(0 \leq y < 10 \rightarrow \neg Q(y)) \Big) \wedge \\
    \Big( (\forall z: \mathbb{Z})(0 \leq z < 10 \rightarrow P(z)) \rightarrow \\
    (\exists v, w: \mathbb{Z})(0 \leq v < 10 \wedge 0 \leq w < 10 \wedge v \neq w \wedge Q(v) \wedge Q(w)) \Big)
    \end{multline*}

\end{enumerate}

\vspace{0.5cm}
\hrule
\vspace{0.5cm}